\section{The actual rational side and moved reader on a given realization}
\label{sec:common-readings-v039}

This chapter constructs the original level-zero rational side and the
original moved reader.  The inputs are one already given geometric
realization and its stored comparison maps.  In particular, existence of
that realization, its primitive rational tests and its comparison maps
is not inferred from the constructions below.  The two source endpoints
are \texttt{exists\_side\_reading\_level\_zero} and
\texttt{read\_of\_moved}.  Together they give \texttt{Read} and
\texttt{RationalReadingWithoutCongruence} on the same data.  The final
\texttt{etaJ} field of \texttt{ReadL} is installed only under its original
agreement condition, proved sufficient in RDR22.

All source references refer to the fixed commit
\texttt{7cbf7a0c32fca1badf05931879a37c3f8a2b03b4}.
The formal proof is assumed correct.  The arguments below reconstruct
the finite linear, integral and character calculations independently;
source theorem names identify the exact contracts rather than replace
their human proofs.

% BEGIN RDR1
\paragraph{RDR1: the fixed data and the order of the two constructions.}
\label{claim:RDR1}
Fix the original
\(B,f,T,w,Z_0,h_X,W_t,z_0,q,o\), where
\(q:\mathrm{RowPk}(w,z_0)\), \(o:q.\mathrm{Over}\), and assume
\(\mathrm{Inputs}(w.D(o.Z))\) and \(o.\mathrm{Fits}\).
Fix one \(\mathfrak m:\mathrm{Moved}(o)\).  Throughout put
\[
 \begin{gathered}
 D=w.D(o.Z),\quad K=B.LL,\quad \sigma=B.\mathrm{embLL},\\
 R=\mathfrak m.z.R,\quad F=\mathfrak m.z.F,\quad
 y=\mathfrak m.y,\quad t=\mathfrak m.t .
 \end{gathered}                                         \tag{RDR.1}
\]
The sharp and base chains are the already stored \(q.p.\mathrm{ch}(1)\)
and \(q.p.\mathrm{ch}(0)\); their maps into the actual multiplicity
spaces are \(t.\mathrm{ch}(1)\) and \(t.\mathrm{ch}(0)\).
The lift comparison is \(t.\mathrm{lf}\).  No one of these is selected
again when a rational side or reader is requested.

Write \(\mathcal R_0(Y)\) for the original rational reading
without congruence on \(R,F,s.p_0,Y\).  The constructions
below prove, in this order,
\[
\begin{split}
 &\exists s:\mathrm{Side0}(o),\quad
   \forall Y:\mathrm{Read}(R,F.\mathrm{lv2},s.p_0),\\
 &\hspace{2em}\mathcal E(Y,y,t)\Longrightarrow
                       \mathrm{Nonempty}\bigl(\mathcal R_0(Y)\bigr),\\
 &\forall s:\mathrm{Side0}(o),\quad
   \exists Y:\mathrm{Read}(R,F.\mathrm{lv2},s.p_0),\quad
                  \mathcal E^{+}(Y,y,t).
\end{split}                                                   \tag{RDR.2}
\]
Here \(\mathcal E\) is exactly the four ties in the first source
endpoint: \(Y.\mathrm{tupT}=y.\mathrm{tupT}\), equality of the sharp
and base \(j_E\) maps, and equality of the \(e_W\) maps.
The stronger \(\mathcal E^{+}\) consists of the twelve equalities
listed in RDR20.  Thus the side \(s\) is chosen before the later
reader \(Y\).  This order is part of the assertion.

\emph{Proof of the data dictionary.}
The record \(\mathrm{Moved}\) supplies \(R,F\), their original
\(\mathrm{Pins}\), the exact \(\mathrm{Holds}(R,F)\) field
\(\mathrm{intersectionNumber\_baseChange}\), and the whole \(\mathrm{PkT}\)
on \(q.p\).  Its remaining fields fix the original trace set, trace
function, core transport, \(\eta\)-function and period data; none is
changed below.  The record \(o.\mathrm{Fits}\) supplies the original
sharp division comparison, the original finite group data and its
central-action comparison for every standard residual isomorphism.
Finally \(q.\mathrm{num}\) and \(q.\mathrm{kap}\) are quantified over
every \(\eta,h_\eta,V,S,N,G\); evaluating them at the already fixed
fields of \(o\) gives their literal instances for \(D\).  This is
evaluation of a supplied common record, not a new existence assertion.
% END RDR1

% BEGIN RDR2
\paragraph{RDR2: an explicit simultaneous localization submodule.}
\label{claim:RDR2}
Let \(V\) be an \(\mathcal O_K\)-module, let \(P_u\) be the
specified prime ideals, and let \(S_u\subset V\) be submodules.
Then
\[
 L=\{x\in V:\ \forall u\ \exists d\in\mathcal O_K,
                     \ d\notin P_u,\quad dx\in S_u\}        \tag{RDR.3}
\]
is a submodule, with exactly the displayed membership criterion.
The index set may be infinite and need not consist of different primes.

\emph{Proof.}
For zero take \(d=1\); a proper prime does not contain one.
For \(x,y\in L\), at a fixed \(u\) choose witnesses \(a,b\).
Primality gives \(ab\notin P_u\), and
\(ab(x+y)=b(ax)+a(by)\in S_u\).  For \(c\in\mathcal O_K\),
the witness \(a\) for \(x\) also works for \(cx\), since
\(a(cx)=c(ax)\).  These prove all module axioms.  The proof
does not assert that \(L\) is finitely generated or spans \(V\).
Those are separate properties in the original notion
\(\mathrm{HLat.Good}\); they are not fields of \(\mathrm{HLat}\)
or of the side constructed here.
% END RDR2

% BEGIN RDR3
\paragraph{RDR3: the exact distinguished Gram frame.}
\label{claim:RDR3}
Let \(e:K\hookrightarrow\mathbb C\), let \(W\) be a finite
dimensional complex Hilbert space, and fix a basis
\((c_i)_{i<n}\).  At \(v_e=\mathrm{mk}(e)\) there is an
invertible frame \(X_e\) such that, for the standard basis of
\(K^n\),
\[
 \bigl\|X_e^{-1}\,v_e(x)\bigr\|_2
                =\left\|\sum_i e(x_i)c_i\right\|.           \tag{RDR.4}
\]
Here \(v_e(x)\) means entrywise application of the source's chosen
embedding representing the infinite place.

\emph{Proof.}
Choose an orthonormal basis \((a_i)\) of \(W\) by finite
Gram--Schmidt, and write \(c_j=\sum_i A_{ij}a_i\).
Since both families are bases, \(A\) is invertible and
\(\|\sum z_ic_i\|=\|Az\|_2\).  If the representative of
\(v_e\) is \(e\), set \(X_e=A^{-1}\).  If it is
\(\overline e\), set \(X_e=\overline A^{-1}\).  In the latter
case \(X_e^{-1}v_e(x)=\overline{Ae(x)}\), which has the same
Euclidean norm.  This proves (RDR.4) in both cases.  It also
checks the direction of the inverse: the source norm uses
\(X_e^{-1}\), not \(X_e\).  For \(n=0\), all spaces and
sums are zero and the empty matrix has determinant one, so
the same assertion holds.
% END RDR3

% BEGIN RDR4
\paragraph{RDR4: the literal Hermitian-lattice record with prescribed local tests.}
\label{claim:RDR4}
Suppose \(X^o(v)\) is an invertible \(n\)-by-\(n\) frame at
every infinite place, and \(S_u\subset W\) are arbitrary sets.
There is an original \(\mathrm{HLat}(K)\), a semilinear map
\(j:RR.V\to_e W\), and an equivalence
\(\epsilon:RR.V\simeq_K K^n\) such that
\[
\begin{gathered}
 jx=\sum_i e((\epsilon x)_i)c_i,\qquad
 \|x\|_{RR,v_e}=\|jx\|,\\
 \{j(RR.b_i)\}_i\text{ is a complex basis of }W,\\
 \mathrm{maj}_{RR,v}(x,x')
   =\mathrm{maj}_{X^o,v}(\epsilon x,\epsilon x')
                                      \quad(v\ne v_e),       \tag{RDR.5}\\
 x\in RR.L\ \Longleftrightarrow\
 \forall u\ \exists d\notin P_u,
     \quad dx\in\operatorname{span}_{\mathcal O_K}
                              \{z:jz\in S_u\}.
\end{gathered}
\]
Every frame of \(RR\) is invertible.

\emph{Proof.}
Take \(RR.V=K^n\), its standard basis, \(\epsilon=1\),
\(j(x)=\sum e(x_i)c_i\), and the submodule in RDR2 with
\(S_u\) there replaced by the span in (RDR.5).  Set
\(RR.X(v_e)=X_e\) from RDR3 and \(RR.X(v)=X^o(v)\)
elsewhere.  Invertibility is immediate in each branch.
The images of the standard basis are exactly \(c_i\).
RDR3 gives the distinguished norm and the identical frame
gives every off-distinguished Hermitian pairing, with the
original conjugation convention in both entries.  The last
line is the definition of \(L\), not a lattice saturation
claim needing a finite generation theorem.  If a dimension
equality reindexes \(\mathrm{Fin}(n)\), reindex both matrix
indices and both coordinate vectors by that same equality;
matrix multiplication and the finite inner-product sum are
unchanged.  This also proves the source's \texttt{majAt\_cast}.
% END RDR4

% BEGIN RDR5
\paragraph{RDR5: original quotient coordinates and invertibility of all row frames.}
\label{claim:RDR5}
For either \(b\in\{0,1\}\), put \(X=R.\mathrm{toCx3}(F)\)
and let \(S_b=X.\mathrm{data}(F.cr,b)\) be the original
quotient-coordinate data.  The map
\(I_b=X.iC(F.cr,b,\mu_0)\) is injective and has range
\(X.EEl_b\); every original frame \(S_b.XE(v)\) is invertible.
These conclusions use the actual \(\mathrm{Pins}\) and
\(\mathrm{Holds}\) from RDR1.

\emph{Proof.}
DI14 identifies the exact coefficient form, rational adjoint cut,
finite-level projector and positive theta operators with TD1--TD9.
In particular the rational cut spans
\(\operatorname{im}\Theta^*\cap V_{0,b}\), and the two
level inclusions are \(\Theta V_{0,b}\subset E_b\) and
\(\Theta^*E_b\subset V_{0,b}\).  The original
\(\mathrm{sumE,posScal,thmB}\) say that \(E_b\) is the
orthogonal sum of its specified type pieces and that
\(\Theta\Theta^*\) is a strictly positive scalar on each
piece.  The argument of TD8 therefore applies to precisely
\(F.cr.\mathrm{toR}\): its composite with \(\Theta j_R\),
which is \(I_b\), is injective and maps onto \(X.EEl_b\).
The zero-dimensional case is included.

Here is the further frame check.  At the distinguished place,
the source Gram matrix is
\[
        H_b=n_P(F,b)^{-1}
              \bigl(\langle I_b e_i,I_b e_j\rangle\bigr)_{ij}.
                                                               \tag{RDR.6}
\]
Since \(n_P(F,b)>0\) and \(I_b\) is injective, this is positive
definite.  Diagonalizing the Hermitian matrix gives an invertible
positive square root; the original frame \(\mathrm{frameOf}(H_b)\)
is its inverse.

We check the quotient metric at every other embedding as well.
Write \(V\) for the actual image-basis matrix of the finite
good-level idempotent, and \(H\) for the original coefficient
form.  The source definitions give
\(H_V=|P|^{-1}V^*HV\).  The action matrices and \(H\) have
entries in \(\mathrm{cmMax}(K)\) by \(\mathrm{Ok3}\).
So does the good-level idempotent: its character and level
character values are sums of roots of unity and roots of unity,
respectively.  Complex conjugation on a cyclotomic field
commutes with every embedding, so their values lying in \(K\)
belong to this conjugation-compatible subfield.  A maximal
independent subset of the columns of that idempotent gives
an image basis over this subfield.  Thus the source's
\(\mathrm{subBasis}\), whose first branch chooses such a
rational basis whenever one exists, has entries in the same
subfield.  Consequently at every \(v\)
\(H_V(v)=|P|^{-1}V(v)^*H(v)V(v)\) is positive definite.

Let \(Z\) be the actual coordinate matrix of the rational
adjoint cut.  At the distinguished embedding its Gram matrix
is invertible, so the original
\(P_q=(Z^*H_VZ)^{-1}Z^*H_V\) satisfies \(P_qZ=1\)
over \(K\), as in SLR2.  That equality remains true after
every embedding, hence \(P_q(v)\) is surjective everywhere.
The exact quotient metric is
\[
 H_Q(v)=\bigl(P_q(v)H_V(v)^{-1}P_q(v)^*\bigr)^{-1}.
\]
The matrix inside the inverse is positive definite: its
quadratic form at \(x\ne0\) is that of \(H_V(v)^{-1}\)
at the nonzero vector \(P_q(v)^*x\).  Thus \(H_Q(v)>0\)
and the original \(X_Q(v)=H_Q(v)^{-1/2}\) is invertible.
This argument requires no unproved conjugation compatibility
of the rational-cut matrix at the other embeddings.

The type projections are joint \emph{generalized}-eigenspace
projections, not finite-group isotypic projections.  Apply
SLR4's polynomial primary-decomposition proof to the actual
commuting split Hecke operators in \(\mathrm{Ok3}\).
TD4--TD7 give their invariant adjoint cut; their commutation
with the finite group makes the same polynomial projections
preserve the good-level image.  The \(\mathrm{block}\)
condition puts the cut in the prescribed block sum.  Each
nonzero primary component there has residual polynomial
\(S_p\) characterized by
\(H_P=S_p(\mathrm{pl}(P))B_P\), where the fixed block
polynomial \(B_P\) has constant term one.  Cancellation
makes the residual label unique; normalize it to one at
unused places.  A nonzero component in the good-level image
also lies in the full good image, so witnesses that this
label occurs in the source's finite \(\mathrm{ixFin}\).
There are only finitely many primary components, hence the
conditional finite-set definition takes its actual finite
branch.  Their inverse images under the injective original
coordinate map therefore give exactly the direct sum of
the \(Q_p\)'s.  The projections \(P_p\) onto these summands
sum to one.  This proves the particular direct-sum and
quotient-frame facts needed here without assuming the
whole well-formedness package or semisimplicity.

The off-distinguished matrix inverted in \(S_b.XE(v)\) is the
positive square root of
\[
       H(v)=\sum_p a_p(v)P_p(v)^*H_Q(v)P_p(v),\qquad
       a_p(v)=\kappa(v)\prod_j\lambda_{j,p}(v)>0.             \tag{RDR.7}
\]
For a nonzero vector, some \(P_p(v)x\) is nonzero because
their sum is \(x\).  Every summand in \(x^*H(v)x\) is
nonnegative and that summand is positive.  Thus \(H(v)\)
is positive definite.  The source's subsequent change of
coordinates is the identity and its splice at \(v_0\)
is the invertible frame already checked.  Hence every
\(S_b.XE(v)\) is invertible, with its original coefficients
\(\kappa\prod_j\lambda_{j,p}\) unchanged.
% END RDR5

% BEGIN RDR6
\paragraph{RDR6: a basis through an involutively semilinear range.}
\label{claim:RDR6}
Let \(\tau:\mathbb C\to\mathbb C\) be an involutive field
automorphism, let \(g:\mathbb C^N\to W\) be additive and
\(\tau\)-semilinear with zero kernel, and let
\(j_E:E\to W\) be injective and complex linear with
\(\operatorname{range}j_E=\operatorname{range}g\).
There is a basis \((c_i)_{i<N}\) of \(E\) such that
\[
                  j_E\left(\sum_i\tau(z_i)c_i\right)=g(z)
                         \qquad(z\in\mathbb C^N).          \tag{RDR.8}
\]

\emph{Proof.}
Let \(c_i\) be the unique preimage under \(j_E\) of
\(g(e_i)\).  Finite additivity and semilinearity give
\(g(z)=\sum_i\tau(z_i)g(e_i)\), proving (RDR.8).
If \(\sum a_ic_i=0\), substitute \(z_i=\tau(a_i)\).
Then \(g(z)=0\), so \(z=0\), and applying \(\tau\)
again gives \(a_i=0\).  For any \(x\in E\), choose \(z\)
with \(j_Ex=g(z)\); injectivity of \(j_E\) and (RDR.8)
give \(x=\sum\tau(z_i)c_i\).  Thus the independent family
spans.  Empty families give the same proof when \(N=0\).
% END RDR6

% BEGIN RDR7
\paragraph{RDR7: the original \texttt{RrT} basis and all its frames.}
\label{claim:RDR7}
For either \(b\), suppose a specified injective linear map
\(j_E:E\to(y.C(F.\mathrm{lv2})).\mathrm{Amb}(0)\) has range
exactly \((y.C(F.\mathrm{lv2})).EEc_b\).  There are
\(n,c,X^o\) satisfying the literal \(\mathrm{RrT}(F,y,b,j_E,n,c,X^o)\).
In particular \(n=S_b.D\), and \(X^o(v)=S_b.XE(v)\)
with the canonical dimension reindexing, at every infinite place.

\emph{Proof.}
Let \(s=\mathrm{twAt}(\sigma,v_0)\), and let \(\tau\)
be identity or complex conjugation according to \(s\).
Define the map into the original tensor-linearized multiplicity space by
\[
 g(z)=\mathrm{TensorLift}\bigl(R.HE.e
                (\mathrm{Tw.out}(I_bz))\bigr).               \tag{RDR.9}
\]
It is additive and \(\tau\)-semilinear.  RDR5 gives
injectivity of \(I_b\); the original \(\mathrm{models.HE}\)
gives injectivity of \(R.HE.e\), and tensor linearization
is an equivalence.  Thus \(g\) has zero kernel.

Its range is exactly \(EEc_b\).  In one direction, RDR5
puts \(I_bz\) in the twisted inverse image of the original
\(R.EEat(R.\mathrm{levR}(F,b))\), so (RDR.9) lies in its
tensor-linearized image.  Conversely every vector of that
image comes from the same \(EEat\).  Its finite principal
level is contained in the level used by \(HE\): the residual
matrix is identity there, hence lies in the selected subgroup,
all its prescribed characters equal their identity values,
and the other level conditions are unchanged.  Thus the
vector has a preimage under \(R.HE.e\); after applying
\(\mathrm{Tw.mk}\), RDR5 puts that preimage in the range
of \(I_b\).  This proves equality of ranges using the
actual level inclusion, with no dimension inference.

Apply RDR6 to this \(g\) and the supplied \(j_E\).
The distinguished embedding dictionary is
\[
                    \tau(v_0(a))=\sigma(a)\quad(a\in K).
                                                               \tag{RDR.10}
\]
It follows from \(v_0=\mathrm{mk}(\sigma)\) and the definition
of \(\mathrm{twAt}\); it also covers the conjugate representative
of that infinite place.  Substituting \(z_i=v_0(x_i)\) into
(RDR.8) gives
\[
 j_E\left(\sum_i\sigma(x_i)c_i\right)
   =\mathrm{TensorLift}\bigl(R.HE.e
         (\mathrm{Tw.out}(I_b(v_0(x))))\bigr),               \tag{RDR.11}
\]
which is the exact coordinate clause of \(\mathrm{RrT}\).
Choose \(X^o=S_b.XE\); RDR5 proves all its determinant
clauses.  Neither a new rational cut nor an isomorphic
replacement of \(E\) has been introduced.
% END RDR7

% BEGIN RDR8
\paragraph{RDR8: the unique standard residual read and its central action.}
\label{claim:RDR8}
Put \(e_0=\mathfrak m.z.cG\circ R.eG\), with composition
understood in the source order \(R.eG.\mathrm{trans}(\mathfrak m.z.cG)\).
Every \(e_R\) satisfying the original \(\mathrm{StdIso}\)
equals \(e_0\).  Pullback of coefficients along this fixed
isomorphism gives an algebra homomorphism on group-algebra centers,
\[
       (\mathrm{cenPull}_{e_0}z)(g)=z(e_0g).                 \tag{RDR.12}
\]
On the original base space define
\(\mathrm{cen}_\mu(z)=
 (q.p.sA(0,\mu)).\mathrm{act}(\mathrm{cenPull}_{e_0}z)\).
This satisfies both central-action fields of \(\mathrm{Side0}\).

\emph{Proof.}
The \(\mathrm{StdIso}\) witnesses for \(e_R\) and the
\(\mathrm{Keep.grp}\) witnesses for \(e_0\) label the same
selected prime over each of the two distinguished rational primes.
The two residue maps from \(\mathbb Z/p\mathbb Z\) to that
residue field agree: every element is represented by an integer,
and any unital ring homomorphism must take that integer to its
same multiple of one.  They are injective because their domain
is a field.  The \(\mathfrak m.z.idG\) matrix-entry equation
therefore forces the two residual matrices to agree at each
coordinate.  This proves \(e_R=e_0\), including both prime
families, without choosing a different finite-group read.

Reindexing a finite convolution sum by the group isomorphism
proves that (RDR.12) is an algebra homomorphism.  It takes a
central element to a central element, since conjugation is
preserved.  The field \(\mathrm{Fits.grpAt.cen}\) says that
the actual \(D.\mathrm{ntypes}.e_0\) is central.  On every
coordinate vector of the spanning \(t.tA(0,\mu).j\) family,
the stored \(\mathrm{FormsTie}\) and matrix law \(E=\mathrm{op}(a)\)
give the same action as (RDR.12), with the coefficient identity
from \(\mathfrak m.z.\mathrm{good}\).  Two complex linear
operators agreeing on that spanning family are equal.  Hence
the stored base \(e\)-operator is exactly
\(\mathrm{cen}_\mu(D.\mathrm{ntypes}.e_0)\); the conditional
zero branch in the source definition is excluded by centrality.

For the universal \(\mathrm{Side0.is}\) clause, first replace
the supplied \(e_R\) by \(e_0\) using uniqueness.  The two
finite group-algebra elements then have identical coefficients
at every group element, hence are equal.  Applying the same
stored action proves that clause for every central element
and every vector.
% END RDR8

% BEGIN RDR9
\paragraph{RDR9: the complete original \texttt{Side0} constructor.}
\label{claim:RDR9}
Given the basis and invertible frames of RDR7 for \(b=0\),
and arbitrary prime-indexed sets \(S_u\) in the original base
\(EE\), there is a \(\mathrm{Side0}(o)\) satisfying every
conclusion of (RDR.5), with \(e=\sigma\).

\emph{Proof.}
Use the exact central action of RDR8 and the exact
\(RR,j,\epsilon\) of RDR4.  These give, respectively,
\(\mathrm{cen,tie,is}\) and
\(\mathrm{RR,frames,jEE,basis,norm}\), all eight fields
of \(\mathrm{Side0}\).  The associated \(s.p_0\) uses
the original base \(\mathrm{Amb,Mult,multZero,multTop,io}\),
the original integral carrier \(\mathrm{Mtilde}\), and
the original \(e\)-operator with its self-adjointness and
idempotence.  It changes only the permitted central reading,
the rational lattice record and its map into the same \(EE\).
Consequently no level, original Haar measure, complex period
or chain map changes in this construction.  This is the
literal source \texttt{exists\_side}, including its arbitrary
sets \(S_u\), not an abstract existence of compatible data.
% END RDR9

% BEGIN RDR10
\paragraph{RDR10: the noncircular two-stage construction of the local lattice.}
\label{claim:RDR10}
There is a side \(s:\mathrm{Side0}(o)\) satisfying the
original predicate \(\mathrm{Good}(\mathfrak m,s)\).
In particular its lattice satisfies its own literal
\(s.p_0.s.\mathrm{RRDef}\).

\emph{Proof.}
We first check the independence that permits the construction.
For two sides built over the same \(q.p\), the recursion
\[
       \mathrm{downTo}_{\mu,0}=1,\qquad
       \mathrm{downTo}_{\mu,k+1}
                   =\mathrm{downTo}_{\mu,k}\circ\mathrm{io}_{\mu,k}
                                                               \tag{RDR.13}
\]
has identical terms, because all \(\mathrm{io}\)'s are copied.
Thus the complete \(\iota_\mu\), which also uses the unchanged
\(\mathrm{multZero}\) and \(\mathrm{multTop}\), is identical.
The space \(V_{0,\mu}\), the kernel cut
\(U_\mu=V_{0,\mu}\cap\ker\iota_\mu\), and
\(Q_\mu=V_{0,\mu}\cap U_\mu^\perp\) are therefore identical
with the same original inner product.  So is the source set
\[
  \mathrm{RRset}_\mu
    =\iota_\mu\bigl(\operatorname{pr}_{Q_\mu}
                          (\mathrm{Mtilde}_\mu\cap V_{0,\mu})\bigr).
                                                               \tag{RDR.14}
\]
In particular it does not depend on \(RR.L\), \(j_{EE}\),
or the new central-action record.  This proves both
\texttt{RRset\_reRow} and \texttt{RRset\_side} by their
actual formulas, not by identifying two unspecified lattices.

Apply RDR7 to the given base map \((t.\mathrm{ch}(0)).j_E\).
Use RDR9 first with the original primes \(D.\mathrm{plEq}(u)\)
and the empty set at each prime, obtaining a temporary side \(s'\).
No fullness of its lattice is needed.  Apply RDR9 a second time,
using the very same basis and frames, now with
\(S_u=(s'.p_0.s).\mathrm{RRset}_{D.\mathrm{muAt}(u)}\).
Call this side \(s\).  Its membership formula is
\[
 x\in s.RR.L\ \Longleftrightarrow\
 \forall u\ \exists d\notin D.\mathrm{plEq}(u),\quad
 dx\in\operatorname{span}_{\mathcal O_K}
 \{z:s.j_{EE}z\in(s'.p_0.s).\mathrm{RRset}_{D.\mathrm{muAt}(u)}\}.
                                                               \tag{RDR.15}
\]
By (RDR.14) replace \(s'\) by \(s\) in this formula.  The
result is exactly \(\mathrm{RRDef}\).  RDR7 supplies the
\(\mathrm{RrT}\) field, and RDR9 supplies the coordinate and
off-distinguished metric fields.  These are all fields of
\(\mathrm{Good}(\mathfrak m,s)\).  This predicate is distinct
from \(\mathrm{HLat.Good}\); no finite generation or full
rank of the newly defined lattice is asserted here.
% END RDR10

% BEGIN RDR11
\paragraph{RDR11: the unique original \(j_W\), with its exact range and normalization.}
\label{claim:RDR11}
On the given \(R,F\), let \(mj:\mathrm{MjTie}(R,F.\mathrm{lv2})\)
be any original lift tie.  For every \(i<D.R.m\) there is a
\(\sigma\)-semilinear map from the original base quotient
coordinates to the original \(D.\mathrm{what}_2.M_j(i+1)\),
uniquely determined by
\[
 mj.eW_i(jW_i(x))
  =(R.FW_i).j(\mathrm{sg0})
                   \left(\sum_k x_k R.bL(F.\mathrm{lv2},0,i)_k\right).
                                                               \tag{RDR.16}
\]
Its source dimension is exactly \(S_0.dW(i)\).

\emph{Proof.}
Write the right side as \(g_i(x)\).  The original
\(\mathrm{sg0}\) is \(\sigma\), since \(R.hv0\) fixes
the same infinite place and its twist convention; hence
\(g_i\) is \(\sigma\)-semilinear.  The source definition
of \(b_C\) gives
\[
 g_i(e_k)=\sqrt{n_P(F,0)}\,b_{C,k},\qquad
       \sqrt{n_P(F,0)}>0.                                  \tag{RDR.17}
\]
Thus the complex span of these coordinate images is the
original space \(R.WS(F,0,i,\mathrm{sg0})\).  The given
\(\mathrm{Holds.aux.levelW}\) identifies this exact space
with the good finite-level subspace, and \(mj.eW\)'s range
field identifies it with \(\operatorname{range}(mj.eW_i)\).
Every \(g_i(x)\) is a complex linear combination of its
coordinate images, so lies in that range.

The other exact field of \(mj\) is
\[
               \|mj.eW_i(z)\|^2=n_P(F,0)\|z\|^2.           \tag{RDR.18}
\]
Since \(n_P(F,0)>0\), this map is injective.  Define
\(jW_i(x)\) to be the unique preimage of \(g_i(x)\).
Applying \(mj.eW_i\) shows that \(jW_i(x+x')\) and
\(jW_i(x)+jW_i(x')\) have the same image; injectivity
gives additivity.  The same argument with a scalar gives
\(jW_i(ax)=\sigma(a)jW_i(x)\).  It proves uniqueness as
well.  The source dimension is the finite rank of the same
\(R.WL(F,0,i)\), precisely the definition of \(S_0.dW(i)\).
This construction supplies \(j_W\) itself; it does not assume
the existence of the \(Z.j_W\) used in DI18.
% END RDR11

% BEGIN RDR12
\paragraph{RDR12: the complete level-zero rational-reading endpoint.}
\label{claim:RDR12}
The first assertion of (RDR.2) holds under exactly the inputs of
RDR1.  Its conclusion has all ten fields of the original
\(\mathrm{RationalReadingWithoutCongruence}\).

\emph{Proof.}
Choose the single \(s\) from RDR10.  Now let \(Y\) be any
reader satisfying the four ties \(\mathcal E\).  The sharp
\(\mathrm{Fits.div}\) field gives an equivalence
\(\epsilon_1:D.\mathrm{div}.RR.V\simeq_KK^{n_1}\),
the exact coordinate expression of its \(j_{EE}\), and
its off-distinguished metric comparison with the stored
\(q.p.X^o\).  The sharp \(\mathrm{RrT}\) is the original
\(t.rr\).  RDR10 gives an equivalence \(\epsilon_0\) and
the same three facts at level zero.  Take
\[
                   eRR=\epsilon_1^{-1},\qquad
                   eRR0=\epsilon_0^{-1}.                   \tag{RDR.19}
\]
The dimensions in the two \(\mathrm{RrT}\)'s are their
actual quotient dimensions.  Substitute (RDR.19) into
their coordinate identities (RDR.11), and use the supplied
equalities of \(Y.I_{\rm sh}.j_E\), \(Y.I_0.j_E\) with
the given comparison maps.  This gives the literal \(jEE\)
and \(jEE0\) fields.  Substitute into the two metric
comparisons and reindex both coordinates and frame indices
as in RDR4.  This gives the literal \(\mathrm{met,met0}\)
fields at every \(v\ne v_0\), with no omitted embedding.

The sharp \(\mathrm{rr}\) field is
\(\mathrm{Fits.div}.2(o.S,o.N,o.G)\); the base \(\mathrm{rr0}\)
field is RDR10.  Finally apply RDR11 to this very \(Y.mj\)
to obtain \(jW,jW\_read\).  Together with the two equivalences,
these are precisely the ten required fields.  The two other
ties in \(\mathcal E\), for \(\mathrm{tupT}\) and \(e_W\),
remain in the original theorem statement, although this
rational-coordinate construction only uses the two \(j_E\)
ties.  No \(\eta_J\) equality has been inferred.
% END RDR12

% BEGIN RDR13
\paragraph{RDR13: the exact finite-level inclusions and real parameters.}
\label{claim:RDR13}
For the original chain associated to \(R,F,y\),
\[
 EEc_b\subset CC_{b,\mu,0},\qquad
 CCP_{b,\mu,j}\subset CC_{b,\mu,j},\qquad
 CCXP_{b,\mu}\subset CCX_{b,\mu}.                           \tag{RDR.20}
\]
The middle inclusion is at the original inner indices and the
last at the original terminal level.  The corresponding
target holomorphic parameters are exactly the stored
\(\mathrm{parW}\), not merely isomorphic parameters.

\emph{Proof.}
A function transforming under a level \((J,\chi)\) transforms
under \((J',\chi')\) whenever \(J'\subset J\) and the two
characters agree on \(J'\).  All its smoothness, automorphy,
cuspidality, equivariance and continuity properties are
unchanged.  The implication therefore holds for each generator
and for its linear span, and also after the original tensor
linearization.

At the principal residual kernel the residual matrix is one.
It belongs to every selected residual subgroup; its determinant
characters equal one.  Thus passing from the source full
level to its selected \(P\)-level, or from the \(P\)-level
back to the principal full level used in (RDR.20), has exactly
the required subgroup inclusion and character restriction.
The other local clauses of \(\mathrm{datY}\) or \(\mathrm{datX}\)
are copied.  At row index zero the two dimension inequalities
force the block size to be the full rank, so the remaining
quotient block has size zero and determinant one.  This
checks the \(\mathrm{dat0}\)-to-\(\mathrm{datY}\) character
identity including that empty block.  These observations
give all three inclusions in (RDR.20).

For inner row indices, the source truncation is
\(\mathrm{parB}(\min(j,m))\); in the used range \(j<m\)
it equals \(\mathrm{parB}(j)\).  At the target, the
distinguished coefficient embedding is \(\mathrm{sg0}=\sigma\),
as in DI16 and (RDR.10).  Substitution in the original ordered
label and weight formulas gives \(\mathrm{SW}(i)=\mathrm{parW}(i)\)
and hence the required holomorphic-function model.  No
integrability assertion is made in this argument.
% END RDR13

% BEGIN RDR14
\paragraph{RDR14: the original row identifications.}
\label{claim:RDR14}
Suppose \(\mathrm{Used0}(D,\mu)\iff y.\mathrm{Us}(\mu)\).
For either original side \(S\), a supplied comparison
\(t_b:\mathrm{ChT}\) constructs a \(\mathrm{RowIds}\)
and a full \(\mathrm{RowSideReal}\), with
\[
 \begin{split}
 (I.j_E x).1&=t_b.j_E x,\\
 (I.j_M(\mu,j)x).1&=t_b.j_M(\mu,j)x,\\
 (I.j_X(\mu)x).1&=t_b.j_X(\mu)(S.\mathrm{multTop}_\mu x),\\
 I.\mathrm{prt}(\pi)&=y.\mathrm{prtE}(F,b,t.\mathrm{rep}(\pi)).
 \end{split}                                                  \tag{RDR.21}
\]
All statements retain their original used-tuple and index guards.

\emph{Proof.}
An injective linear map with specified image is an equivalence
onto that image: its inverse sends a vector to its unique
preimage, and uniqueness proves linearity.  Apply this to
\(t_b.j_E\), to \(t_b.j_M(\mu,j)\) for
\(1\le j\le m-1\), and to \(t_b.j_X(\mu)\).
For the last compose with the original isometry
\(S.\mathrm{multTop}_\mu\).  Define the parts by the last
line of (RDR.21).  This proves the exact maps and their ranges.

RDR13 gives the parameter and level fields of
\(\mathrm{RowSideReal}\).  Its three Hilbert pairing fields
are the corresponding supplied \(j_E,j_M,j_X\) pairing
equalities.  In the terminal one, first use the original
isometry \(\mathrm{multTop}\) to preserve the inner product.
For the \(\mathrm{io0}\) equation, (RDR.20) puts the image
of \(j_M(\mu,1)\) in the required full level, so the exact
\(t_b.\mathrm{io0}\) equation applies.  For every inner
\(j\), do the same with \(j_M(\mu,j+1)\) and
\(t_b.\mathrm{io}\).  At the top use \(CCXP\subset CCX\)
and \(t_b.\mathrm{ioTop}\); the two appearances of
\(m-1+1\) and \(m\) are the canonical index equality.
The separate source clause for a one-step chain is vacuous,
since the original rank datum has \(m\ge2\).

For each orthonormal basis vector with tag \(\pi\), membership
in the span of the vectors with that tag is immediate.
The exact \(t_b.\mathrm{ft}\) identity then puts its image
in the specified part.  The pairwise-disjoint-parts field is
\(t_b.\mathrm{disj}\).  This verifies every field of
\(\mathrm{RowSideReal}\).  Its pairing equalities are part of
the given \(\mathrm{ChT}\); this constructor does not attempt
to manufacture them from an arbitrary \(\mathrm{MMap}\).
In particular no stronger decay theorem is used here.
% END RDR14

% BEGIN RDR15
\paragraph{RDR15: all original lift identifications and the exact \(M_j\) tie.}
\label{claim:RDR15}
Under the same used-tuple equivalence and the original
\(\mathrm{LamDef}(D)\), the given \(t_L:\mathrm{LfT}\)
constructs \(I_L:\mathrm{LiftIds}\),
\(mj:\mathrm{MjTie}(R,F.\mathrm{lv2})\), and
\(\mathrm{LiftTie}(I_L)\).  Their maps satisfy
\[
 (I_L.j_{Mj,i}x).1=t_L.j_{Mj,i}x,\quad
 (I_L.j_{Mt,\mu,i}x).1=t_L.j_{Mt,\mu,i}x,\quad
 mj.eW_i=t_L.eW_i.                                          \tag{RDR.22}
\]

\emph{Proof.}
Make the first two maps equivalences onto their specified
ranges by the same unique-preimage construction as RDR14.
Set \(I_L.\mathrm{prtW}(i,\pi)=y.\mathrm{prtW}
(F,i,t.\mathrm{rep}(\pi))\).  For \(mj\), use exactly
\(t_L.eW\), its given range, its norm equation (RDR.18),
and its given integral membership equivalence
\[
 x\in M_{j,i}\ \Longleftrightarrow\
 \exists v\in (R.FW_i).MK\cap (R.FW_i).W_{\rm good},
                \quad t_L.eW_i(x)=(R.FW_i).j(\mathrm{sg0})v.
                                                               \tag{RDR.23}
\]
This retains both the original coefficient-field norm lattice
and the exact inverse-character good subspace.

RDR13 proves the holomorphic-parameter field.  The two Hilbert
pairing fields of \(\mathrm{LiftTie}\) are the given
\(t_L.jMj\_ip\) and \(t_L.jMt\_ip\).  Given the actual
theta equation on the specified level and the required good
target value, its theta field is precisely \(t_L.\mathrm{theta}\),
with the used-tuple guard transported unchanged.  Its part
and disjointness fields are \(t_L.\mathrm{isot,disj}\).

For the last scalar field, take a nonzero specified type
subspace on which the original endomorphism acts by the real
scalar \(c\).  Choose any basis of that subspace.  Its
matrix is \(cI\), so its complex trace is \(c\) times
the positive dimension; taking real parts and dividing by
that dimension gives exactly \(c\).  The given
\(\mathrm{LamDef}\) identifies the original \(\lambda\)
with this normalized trace.  This proves \(\mathrm{lamP}\)
with the exact nonzero guard, completing \(\mathrm{LiftTie}\).
Neither (RDR.23) nor the original theta equation is inferred
from equality of dimensions; both remain supplied fields of
the common \(\mathrm{LfT}\).
% END RDR15

% BEGIN RDR16
\paragraph{RDR16: the exact quotient of types and occurrence in the two bases.}
\label{claim:RDR16}
The given \(\mathrm{Reps}\) produces an equivalence
\(e_{FT}\) from the original finite type index set to the
original \(\mathrm{FTR}\) quotient.  Its source representative
is isomorphic to \(t.\mathrm{rep}(\pi)\), and this change
of representative preserves every original sharp, base and
lift part.  Each \(\pi\) occurs among the tags of at least
one of the two given row bases.

\emph{Proof.}
Send \(\pi\) to the isomorphism class of
\(t.\mathrm{rep}(\pi)\).  It belongs to the quotient by
\(\mathrm{Reps.mem}\).  Equality of classes is an isomorphism,
so \(\mathrm{Reps.inj}\) makes this map injective.
The supplied completeness clause \(\mathrm{Reps.all}\)
shows that every quotient representative is isomorphic to
one in this list; hence the map is onto.  Taking its inverse
gives the equivalence, with no new representation catalogue.

If \(\rho\simeq\rho'\), compose each intertwiner with that
isomorphism or its inverse.  The resulting intertwiners have
the same images.  Hence their span, the actual isotypic part,
is unchanged.  For the row space apply the same argument to
the conjugate representation; the conjugate of the given
isomorphism is an isomorphism.  This proves the three required
part equalities.

By \(\mathrm{Reps.mem}\), the part for this type is nonzero
in at least one row level \(b\).  Its tensor-linearized part
is still nonzero because tensor linearization is injective,
and it lies in \(EEc_b=\operatorname{range}(t_b.j_E)\).
The \(t_b.\mathrm{ft}\) field says that its inverse image
is spanned by exactly the row basis vectors tagged \(\pi\).
If none were tagged \(\pi\), that inverse image would be
zero.  Every vector of the part has a preimage under \(j_E\),
so the part itself would be zero, a contradiction.  This
proves the required two-level disjunction without requiring
that the type occur in both levels.
% END RDR16

% BEGIN RDR17
\paragraph{RDR17: the trivial base idele character and the period quotient.}
\label{claim:RDR17}
On the same original character model,
\(R.\mathrm{toLevCh}.\xi(\mu_0)=1\).  Consequently
\[
 K_c.\xi(\mu)=R.\mathrm{toLevCh}.\xi(y.\mathrm{tupT}(\mu))
                                                               \tag{RDR.24}
\]
gives the literal \(\mathrm{ColTieR}\), using the original
\(\mathrm{BuDefQ}\) field.

\emph{Proof.}
We first prove the exact character uniqueness required at the
base tuple.  If \(\mathrm{ChrEq}(c,1)\), then \(c\) is
continuous, is one on principal ideles, and is one on every
single finite-place idele of a nonzero element of \(E\).
The actual finite tensor algebra has components
\(\prod_{P\mid v}E_P\), with its restricted-product topology.
The explicit tensor/Chinese-remainder and integral-lattice
argument in JIO14 identifies these with the original
\(AE\) components.  For one place, \(E^\times\) is dense
in \(E_P^\times\), by the definition of completion and the
openness of its nonzero locus.  Continuity therefore makes
\(c\) one on every single-place unit.  Products of these
are dense in the finite idele group: a basic restricted-product
neighborhood prescribes only finitely many coordinates and
allows the full integral unit groups elsewhere, where identity
is permitted.  Thus \(c\) is one on all finite ideles.

For a principal idele, its finite value is now one, so its
real value is one too.  The original CM tensor isomorphism
identifies the real algebra with a product of copies of
\(\mathbb C\).  A rational basis of \(E\) is a real basis
of this real algebra; rational coefficients are dense in
real coefficients.  The nonzero locus is open, so the image
of \(E^\times\) is dense in the real unit group.  Continuity
makes \(c\) one there as well.  The real/finite factorization
proves \(c=1\) on the actual \(AE^\times\).

The given \(\mathrm{BaseW.lev.xi}\) at the base tuple,
and \(\mathrm{xi\_base}\), give precisely
\(\mathrm{ChrEq}(R.\mathrm{toLevCh}.\xi(\mu_0),1)\).
This proves the assertion.  Since \(y.\mathrm{tup}(\mu_0)\)
is the original base tuple and \(\mathrm{tupT}\circ\mathrm{tup}=1\),
(RDR.24) has \(K_c.\xi(\mathrm{base})=1\).
Expanding \(\mathrm{mu}(t)_j=\mathrm{mu0}_j\xi(t)\)
then gives the exact character-tuple equation of
\(\mathrm{ColTieR}\), with no missing base factor.

It remains to verify its period clause.  For a word \(u\),
let \(F_0\in V_0(\mu_0)\), let
\(n=\operatorname{pr}_{U}F_0\), with
\(U=V_0(\mu_0)\cap\ker\iota\).  Then
\(F_0-n\in Q=V_0\cap U^\perp\),
\(\iota(F_0-n)=\iota F_0\), and the canonical section of
\(\iota|_Q\) returns \(F_0-n\).  The source
\(\mathrm{BuDefQ}\) evaluates \(Bu\) on exactly that
section.  Whenever the stated period functional vanishes
on \(U\), its additivity gives
\[
 \mathrm{perCol}_u(\varphi,F_0)
    =\mathrm{perCol}_u(\varphi,F_0-n)
    =Bu_u(\varphi,\iota F_0).                              \tag{RDR.25}
\]
This is the full universal period clause, with its explicit
kernel-vanishing premise retained.
% END RDR17

% BEGIN RDR18
\paragraph{RDR18: the exact used tuples, Hilbert lift and unused-chain vanishing.}
\label{claim:RDR18}
The tuple fields of the moved reader can be taken literally
from \(y\).  They satisfy
\[
 D.\mathrm{muAt}(v)
      =y.\mathrm{tup}(R.\mu(D.\mathrm{plEq}(v))),\qquad
 \mathrm{Used0}(D,\mu)\ \Longleftrightarrow\ y.\mathrm{Us}(\mu).
                                                               \tag{RDR.26}
\]
There is a linear map \(e_E\) into the original \(R.HE.H\)
whose tensor-linearized \(HE\) image is the given sharp
\(j_E\).  For every unused tuple, the original sharp
\(\iota_\mu\) is zero.

\emph{Proof.}
For the first equality, use the given \(\mathfrak m.\mathrm{tup}\)
to write \(y.\mathrm{tup}\) as the original \(\mathrm{ixOf}\).
The map \(w.I.\xi\) is injective by \(w.\mathrm{ok}\),
so it suffices to compare its values.  If \(v\in W_t\),
\(q.\mathrm{mu.on}\) replaces \(\mathrm{muAt}\) by
\(\mathrm{muTw}\), and \(q.\mathrm{mu.tw}\), \(w.C\)'s
fixed frame equality and \(\mathfrak m.z.\mu\) identify both
values with \(T.\mathrm{xiTw}(f,v)\).  If \(v\notin W_t\),
\(q.\mathrm{mu.off}\) and \(w.l.\mathrm{xi0}\) make
both values one.  These are all cases.  The definitions of
\(\mathrm{Used0}\) and \(y.\mathrm{Us}\) are the base tuple
or a tuple occurring in (RDR.26), so they agree.  In either
case \(y.\mathrm{tup}(y.\mathrm{tupT}(\mu))=\mu\):
use the base equality and left inverse in the base case,
and that same left inverse in the occurrence case.

Let \(k=\mathrm{TensorLift}\circ R.HE.e\).  It is injective.
Every value of \((t.\mathrm{ch}(1)).j_E\) lies in \(EEc_1\),
and the principal-level inclusion proved in RDR7 puts it in
\(\operatorname{range}k\).  Define \(e_E\) by the unique
inverse of \(k\) on that range.  Its additivity and linearity
follow by applying the injective \(k\); by construction
\[
          (t.\mathrm{ch}(1)).j_E(x)
              =\mathrm{TensorLift}(R.HE.e(e_E x)).          \tag{RDR.27}
\]
This uses the original Hilbert model rather than an arbitrary
space of the same dimension.

Finally \(t.\mathrm{ch}(1).\mathrm{junk}\) sets every
\(\mathrm{io}_{\mu,j}\) to zero for \(\mu\) outside
\(y.\mathrm{Us}\).  For each positive \(k\), the last
factor in the recursion (RDR.13) is zero, so
\(\mathrm{downTo}_{\mu,k}=0\).  Since \(m>0\), the
stored composite \(\iota_\mu\), including its end
equivalences, is zero.  The used-tuple equivalence makes
this the exact \(\mathrm{Read.junk}\) assertion.
% END RDR18

% BEGIN RDR19
\paragraph{RDR19: coordinate equality transports forms, integral sets and centers.}
\label{claim:RDR19}
The equalities \(t.hA(b,t_0)\), \(t.hL(i,t_0)\)
give canonical reindexings of the original finite coordinates.
They preserve both \(\mathrm{FormsTie}\) and \(\mathrm{OLatTie}\),
and are the only coordinate changes used in the moved reader.
The original central-action fields hold for every \(c_G\)
satisfying \(\mathrm{Keep}\), on both row sides.

\emph{Proof.}
For a matrix-record equality \(M=M'\), the dimension
equality gives \(J:K^{\mathrm{Fin}(M'.n)}\simeq_K
K^{\mathrm{Fin}(M.n)}\).  It is just reindexing by that
equality: after identifying the equal natural numbers,
\(Jx=x\) entry by entry.  The transported semilinear map is
the original \(j\circ J\).  Substituting this identity in
the two pairing/operator equations of \(\mathrm{FormsTie}\)
gives exactly the transported forms.  Substituting it in
each direction of the integral-generator membership
equivalence gives exactly \(\mathrm{OLatTie}\).  Neither
the integral coefficient ring nor its image in \(K\)
changes.  This proves the source coordinate-cast lemmas,
including their equality on individual vectors.

Fix now any \(c_G\) with \(\mathrm{Keep}\).  Its
\(\mathrm{grp}\) field, the equality of its original
selected-prime set with \((B.\mathrm{primes}(w.px)).vGs\),
and \(\mathfrak m.z.idG\) show that
\(R.eG.\mathrm{trans}(c_G)\) satisfies \(\mathrm{StdIso}\).
For a central element \(z\) and a rational group-algebra
element \(a\) with the prescribed coefficient equations,
the original \(\mathrm{Tie1.cen}\) gives
\[
 \mathrm{act}(\mathrm{cenPull}_{e_R}z)(jx)
                           =j\bigl(\mathrm{op}(a)x\bigr).  \tag{RDR.28}
\]
On the base side, \(s.\mathrm{is}\) identifies its central
action with this very pullback for every standard read.
On the sharp side, \(\mathrm{Fits.cen}(e_R)\) gives the
same identity for \(o.N\)'s original central action.
Apply (RDR.28) to the transported row ties just proved.
This is exactly \(\mathrm{CenTie}\) for the two original
row-coordinate maps, universally over \(c_G\), not just
for the distinguished \(\mathfrak m.z.cG\).
% END RDR19

% BEGIN RDR20
\paragraph{RDR20: the complete original \texttt{read\_of\_moved}.}
\label{claim:RDR20}
For every \(s:\mathrm{Side0}(o)\), there exists a
\(Y:\mathrm{Read}(R,F.\mathrm{lv2},s.p_0)\) with the
following twelve literal ties:
\begin{enumerate}
\item \(Y.\mathrm{tup}=y.\mathrm{tup}\) and
      \(Y.\mathrm{tupT}=y.\mathrm{tupT}\).
\item \(Y.PY=y.PY\) and \(Y.PW=y.PW\).
\item \(Y.j_A(t_0)\) and \(Y.j_{A,0}(t_0)\) are precisely
      the given sharp and base \(t.tA\) maps transported by
      \(t.hA\), for every original \(t_0\).
\item \(Y.eMt(i,t_0)x=x\) under the canonical dimension
      equality, for every \(i,t_0,x\).
\item The selected source representative of \(Y.eFT(\pi)\)
      is isomorphic to \(t.\mathrm{rep}(\pi)\), for every \(\pi\).
\item \((Y.I_{\rm sh}.j_E x).1=(t.\mathrm{ch}(1)).j_E x\)
      and \((Y.I_0.j_E x).1=(t.\mathrm{ch}(0)).j_E x\).
\item \(Y.mj.eW_i(x)=t.\mathrm{lf}.eW_i(x)\), and
      \((Y.I_L.j_{Mt,\mu,i}x).1=t.\mathrm{lf}.j_{Mt,\mu,i}(x)\)
      for every original used tuple, index and vector.
\end{enumerate}
Counting the two equations in items 1, 2, 3, 6 and 7 separately
gives the twelve conjuncts of the source theorem.

\emph{Proof.}
Set \(h_B=\mathfrak m.z.\mathrm{pins.char.base}\) and
\(h_N=q.\mathrm{num}(o.\eta,o.h_\eta,o.V,o.S,o.N,o.G)\).
RDR18 supplies the exact used-tuple equivalence.  Apply
RDR14 to the given sharp chain and to the base chain of
\(s.p_0\), obtaining \(I_{\rm sh},I_0\); both use the
given \(t.\mathrm{ch}\) maps.  Apply RDR15 to the given
\(t.\mathrm{lf}\) and \(h_N.\mathrm{lam}\), obtaining
\(I_L,mj\).  Use \(e_{FT}\) of RDR16, \(K_c\) of RDR17,
and \(e_E\) of RDR18.

For clarity we verify all remaining \(\mathrm{Read}\) fields,
including those not present among the twelve displayed ties.
The tuple maps, their left inverse and base equality, the
positive chain length, top multiplicity property and
injectivity of \(\mathrm{extV}\) are copied from \(y\).
RDR18 gives \(\mathrm{tupTie,muAt}\).  The fields
\(PY,PW,\iota_W,\mathrm{isomFin,frameMatch},h_F,\mathrm{stdPsi}\)
are also the same fields of \(y\).  Thus all finite and real
frames refer to the unchanged realization.

The four components of \(\mathrm{tied}\) are the two
\(\mathrm{RowSideReal}\)'s of RDR14, \(\mathrm{LiftTie}\)
of RDR15 and \(\mathrm{ColTieR}\) of RDR17.  RDR16's
representative invariance gives \(\mathrm{prtE,prtE0,prtWE}\).
The \(\xi\) field follows directly from (RDR.24).
Define \(j_A,j_{A,0}\) using the two \(t.hA\) coordinate
equalities, and \(eMt\) by the \(t.hL\) reindexing.
RDR19 proves \(\mathrm{formsA,latA,formsAtBottomLevel,latA0}\)
and \(\mathrm{formsL,latL}\), the latter using the original
\(q.p.tL\).  This retains the full integral sets rather
than only their complex spans.

For \(\mathrm{topRead}\), substitute (RDR.21).  Its left
side becomes
\[
 (t.\mathrm{ch}(1)).j_X\!\left(
   \mathrm{multTop}\bigl(\mathrm{multTop}^{-1}(j_Ax)\bigr)\right)
           =(t.\mathrm{ch}(1)).j_X(j_Ax).                   \tag{RDR.29}
\]
The supplied \(\mathrm{ChT.top}\) now gives exactly
\(R.cX.\mathrm{toMap}(R.prX)((R.HX).e(R.jX(\mathrm{rowForm}(x))))\).
The same cancellation proves \(\mathrm{topRead0}\).
For \(\mathrm{mtRead}\), insert (RDR.22), then the coordinate
identity of RDR19, and then the given \(\mathrm{LfT.read}\).
The resulting right side is the unchanged
\(R.\mathrm{extV}(R.jV(\mathrm{liftForm}(x)))\).
For \(\mathrm{mjRead}\), use \(\mathrm{LfT.jMj\_read}\)
and \(mj.eW=t.\mathrm{lf}.eW\) in (RDR.22).

The \(\mathrm{stable,stable0}\) fields are the original
\(t.tA(1,\mu).\mathrm{stable}\) and
\(t.tA(0,\mu).\mathrm{stable}\), for every tuple.
RDR18 proves \(e_E\), its read and \(\mathrm{junk}\).
The full dimension field \(\mathrm{dall}\) is the exact
\(t.\mathrm{dall}\) equality.  RDR19 proves both universally
quantified central fields.  The remaining numerical and
period fields are
\[
 \mathrm{lam}=h_N.\mathrm{lam},\quad
 \mathrm{kap}=q.\mathrm{kap}(o.\eta,o.h_\eta,o.V,o.S,o.N,o.G),
 \quad \mathrm{bu}=h_N.\mathrm{bu},\quad
 \mathrm{ze}=h_N.\mathrm{ze}.                               \tag{RDR.30}
\]
The occurrence field is the two-level alternative in RDR16.
In particular the universal scalar-uniqueness guard in
\(\mathrm{ze}\) and the original \(\mathrm{BuDefQ}\) are
retained in full.

This lists every field in \(\mathrm{Read}\).  By their
definitions, the tuple, frame and row-map equalities in
the statement hold literally; the dimension reindexing
is identity after transporting equal indices.  RDR14--RDR16
give the remaining map and representative ties.  Thus all
twelve source conjuncts hold simultaneously for this one
reader, for all tuples and indices, and for the arbitrary
given side \(s\).
% END RDR20

% BEGIN RDR21
\paragraph{RDR21: the same original \texttt{ReadPk} and all five pins.}
\label{claim:RDR21}
Under RDR1's original hypotheses there exist
\(s:\mathrm{Side0}(o)\) and \(r:\mathrm{ReadPk}(\mathfrak m,s)\)
such that the literal \(r.\mathrm{Pin}\) holds.  The
\(Y\) and rational reading in this packet use the same
\(R,F,y,t\) and the same side \(s\).

\emph{Proof.}
Choose \(s\) first by RDR12.  Apply RDR20 to that side,
obtaining its single \(Y\) and all twelve ties.
Four of these satisfy exactly the antecedent \(\mathcal E\)
in RDR12, which therefore supplies a rational reading
\(Z\) on this same \(Y\).  Set \(r.Y=Y,r.Z=Z\).
The tuple, frame, row-coordinate and \(eMt\) fields of
\(\mathrm{ReadPk}\) are the corresponding ties of RDR20.
Its separate \(\mathrm{Pin}\) record consists of the
representative isomorphism, the sharp and base \(j_E\)
equalities, the \(e_W\) equality and the \(j_{Mt}\)
equality.  They are the five remaining ties of RDR20.
Thus every dependent type in \(r\) refers to the one
previously selected \(s\) and \(Y\).  The apparent
choice cycle between a rational side and a reader has
been removed by (RDR.14), not by selecting separate
compatible witnesses afterward.
% END RDR21

% BEGIN RDR22
\paragraph{RDR22: exactly when the full \texttt{ReadL} is installed.}
\label{claim:RDR22}
Add the original moved-row agreement
\[
           o.\eta(j,v)=q.\mathrm{eta}(j,v)
                       \quad(1\le j\le D.R.m, \text{all }v).
                                                               \tag{RDR.31}
\]
Then the rational reading of RDR21 extends to the original
\(\mathrm{ReadL}(R,F,s.p_0,Y)\).  More generally, any
\(\mathrm{RationalReadingWithoutCongruence}\) extends
whenever its original \(\eta_J\) agrees with
\(\mathrm{etaOf}(R,F)\) in that same range.

\emph{Proof.}
All fields except \(\eta_J\) are copied from \(r.Z\).
By the definition of \(q.p.\mathrm{what}_2\) and the
assembly of \(o.Z\), the stored \(D.\mathrm{what}_2.\eta_J(j,v)\)
is exactly \(o.\eta(j,v)\).  The \(\mathrm{Moved.eta}\)
field says
\(\mathrm{etaOf}(R,F)(j,v)=q.\mathrm{eta}(j,v)\).
Combining it with (RDR.31) gives the required equality
with the original stored \(\eta_J\).
For \(j=i+1\), where \(i<m\), the range guard
\(1\le j\le m\) holds.  Write the original geometric value as
\[
 \begin{aligned}
 \eta_i^{\rm geom}(v)=A_3.A.\mathrm{fam}.\eta_J\bigl(
   &F.P_0,F.\psi_0,A_3.\mathrm{fam3}.\mathrm{toSub2}(X.\mathrm{zs}),\\
   &i,R.\mu(D.\mathrm{plEq}(v)),D.\mathrm{plEq}(v)\bigr).
 \end{aligned}
\]
Expanding \(\mathrm{etaOf}\) at this index gives exactly
\[
 D.\mathrm{what}_2.\eta_J(i+1,v)=
 \begin{cases}
 \eta_i^{\rm geom}(v),    &D.\mathrm{plEq}(v)\in R.\mathrm{twPl},\\
 0,                       &D.\mathrm{plEq}(v)\notin R.\mathrm{twPl}.
 \end{cases}                                                   \tag{RDR.32}
\]
Here \(X=R.\mathrm{toCx3}(F)\), exactly as in the
definition of \(\mathrm{ReadL}\).  This proves its last
field.  With the unchanged \(\mathfrak m.z\), its given
\(\mathrm{Holds}\), \(Y\), and the original \(\mathrm{ixOf}\)
tuple equality, this is also the reading part of the
source's \(\mathrm{RowsAre}\) installation.  For an
arbitrary \(o.\eta\), (RDR.31) is not automatic; RDR21's
unconditional endpoint remains the rational reading
without congruence.
% END RDR22

\subsection*{Remaining scope on the common construction}

The chapter closes the actual finite rational-side and moved-reader
constructors, including the noncircular local lattice, its distinguished
Gram frame, every off-distinguished pairing, the original target
coordinate map, all reader fields and all pins.  It adds no new
integrability or abstract intertwiner premise.  Its supplied
\(\mathrm{ChT}\) and \(\mathrm{LfT}\) already contain the precise
pairing and theta comparisons consumed in RDR14--RDR15.

The original producers remain substantial.  One must still produce
the common \(z_0,q,o,\mathfrak m\), including the same geometric
realization, the full \(\mathrm{Pins,Holds}\), the original
\(\mathrm{PkT}\) with its two \(\mathrm{ChT}\)'s,
\(\mathrm{LfT}\), representation catalogue and coordinate ties,
and the sharp division, group and central data of
\(\mathrm{Fits}\).  RDR5 consumes the existing primitive row tests,
positive operators and decomposition on that data; it does not
construct them.  RDR11 similarly consumes the exact
\(\mathrm{levelW}\) and \(\mathrm{MjTie}\) range, norm and integral
data.  The numerical \(\mathrm{NumOK,KapDef}\) and period records
are evaluated on their common supplied stages, with their full guards.
The extra \(\eta_J\) agreement is retained precisely in RDR22.

No \(\mathrm{ReadC}\), \(\mathrm{ratG}\), \(\mathrm{MLat}\),
stronger \(\mathrm{DivRows3}\), new common selector, or common
threshold is inferred.  These are pointwise constructors for every
given original moved realization.  They can be used inside the
existing common-family quantifiers without changing their order,
but do not supply the outer existence of that family or the
original universal \(\mathrm{NEb}\) data.
